%!TEX root=../../main.tex

\begin{algorithm}[tb]
	\caption{Approximate ReLU Pruning}
	\label{alg:pruning-solutions-nn}
	\begin{algorithmic}
		\STATE {\bfseries Input:} data matrix \( Z \), weights \( W^{(1)}, W^{(2)} \),
		score function \( s \).
		\STATE \( m \gets \abs{\act(W^{(1)})} \)
        \STATE \( (\leftindex_0{W}^{(1)}, \leftindex_0{W}^{(2)}) \gets (W^{(1)}, W^{(2)}) \).
        \STATE \( q^0_i \gets (X \leftindex_0{W}^{(1)}_{\cdot, i})_+ \leftindex_0{W}^{(2)}_{i} \)
		\FOR {\( k = 0 \text{ to } m -1 \)}
		\STATE \( j^k = \argmin_{i \in \act(\leftindex_k{W}^{(1)}_{\cdot, i})} (\leftindex_k{W}^{(1)}_{\cdot, i}) \)
		\STATE \( \beta^k = \argmin_{\beta} \norm{ \sum_{i \neq j^k} \beta_{i} q^k_i - q_{j^k}^k}_2^2 \)
		\STATE \( i^k \gets \argmax_{i} \cbr{|\beta_i| : i \in \act(\leftindex_k{W}^{(1)}_{\cdot, i})}  \)
		\STATE \( t^k \gets 1/|\beta_{i^k}| \)
        \STATE \( (\leftindex_{k+1}{W}^{(1)}_{\cdot, i}, \leftindex_{k+1}{W}^{(2)}_{i} \gets (\leftindex_k{W}^{(1)}_{\cdot, i}, \leftindex_{k}{W}^{(2)}_{i})
		\cdot (1 - t^k \beta_i)^{1/2} \)
		\STATE \( q^{k+1}_i \gets q_i^{k} \cdot (1 - t^k \beta_i) \)
		\ENDFOR
        \STATE {\bfseries Output:} final weights \( \leftindex_{k}{W}^{(1)}, \leftindex_{k}{W}^{(2)} \)
	\end{algorithmic}
\end{algorithm}



Now we specialize our results for CGL to two-layer neural networks
with ReLU or gated ReLU activations.
We state and prove our results for ReLU networks, but they are easily
adapted to gated ReLUs.
We start by interpreting conditions for uniqueness in the context
of non-convex ReLU models and then move on to discussing optimal pruning
for ReLU networks and continuity properties of the solution function.
Proofs are deferred to \cref{app:specializations}.
%This discussion cannot cover the full range of consequences of our theory,
%so we focus on major topics.

\textbf{Optimal Sets and Uniqueness}:
Combining the mapping between solutions for the convex reformulation
and the original non-convex training problem (\cref{eq:convex-to-relu})
and \cref{prop:sol-fn-cgl}
immediately allows us to characterize the solution set
for the full ReLU problem:
\begin{restatable}{theorem}{reluSolFn}\label{cor:relu-sol-fn}
    Suppose \( m \geq m^* \) and \( \tilde D = \calD_Z \) (no sub-sampling),
    with \( p = |\tilde \calD_Z| \).
    Then the optimal set for the ReLU problem up to permutation/splitting
    symmetries is given as follows:
    \begin{equation}\label{eq:relu-sol-fn}
        \begin{aligned}
            \hspace{-0.5em} \calO_\lambda
            = & \,
            \big\{
            (W^{(1)},  W^{(2)}) :
            \, f_{W^{(1)}, W^{(2)}}(Z)  =  \hat y,
            W^{(1)}_{\cdot, i} = (\sfrac{\alpha_{i}}{\lambda})^{\sfrac{1}{2}} v_i, \\
              & W^{(2)}_{i} = \xi_i (\alpha_i \lambda)^{\sfrac{1}{2}},
            \alpha_i \geq 0, \, i \in [2p] \setminus \calS_\lambda \Rightarrow \alpha_i = 0
            \big\}.
        \end{aligned}
    \end{equation}
\end{restatable}
Eq.~\eqref{eq:relu-sol-fn} abuses notation by using $\calS_\lambda$ as
a subset of the neuron indices \( \cbr{1, \ldots 2p} \), where
\( \cbr{1, \ldots p} \) index positive neurons for which \( \xi_i = +1 \)
(corresponding to blocks \( D_i Z \))
and \( \cbr{p+1, \ldots, 2p} \) index the negative neurons (corresponding to \(
- D_i Z \)), for which \( \xi_i = -1 \).
\cref{fig:solution-set} plots the first feature of three neurons
as they vary over this solution set.
We see that the mapping from convex to non-convex parameterization transforms
the convex polytope of solutions into a non-convex, curved manifold.
The optimal sets for the convex and non-convex problems both live in a
low-dimensional space.

Choosing a sub-sampled set of patterns
\( \tilde \calD \subset \calD_Z \) corresponds to finding a stationary
point of the non-convex training problem \citep{wang2022hidden}.
Using this fact with \cref{cor:relu-sol-fn} finally justifies
description of all stationary points of the ReLU
problem given in the introduction.
\begin{restatable}{proposition}{stationaryPoints}\label{prop:stationary-points}
    The set of stationary points of two-layer ReLU networks up to
    permutation/splitting symmetries is
    \begin{equation*}
        \begin{aligned}
            \calC_\lambda = \big\{
             & (W^{(1)},  W^{(2)}) : \tilde \calD  \subseteq  \calD_Z,
            \, f_{W^{(1)}, W^{(2)}}(Z)  =  \hat y_{\tilde \calD},                                         \\
             & W^{(1)}_{\cdot, i} = (\sfrac{\alpha_{i}}{\lambda})^{\sfrac{1}{2}} v_i(\tilde \calD), \,
            W^{(2)}_{i} = \xi_i(\alpha_i \lambda)^{\sfrac{1}{2}},                                         \\
             & \alpha_i \geq 0, \, i \in [2 |\tilde \calD|] \setminus \calS_\lambda \implies \alpha_i = 0
            \big\},
        \end{aligned}
    \end{equation*}
    where \( \tilde \calD \) are sub-sampled activation patterns,
    \( \hat y_{\tilde \calD} \) is the optimal model fit using
    those patterns, and
    \( v_i(\tilde \calD) = \ci(\tilde \calD) - \Ki \ri(\tilde \calD) \)
    is determined by the fit and the dual parameters.
\end{restatable}

We note that since deeper networks are also related
to CGL through convex reformulations \citep{ergen2021beyond},
\cref{prop:stationary-points} may also be applied beyond two layers.

Recall that the solution set for a two-layer ReLU network
is typically not unique due to model symmetries.
However, if the convex solution is unique, then
the non-convex ReLU training problem is p-unique.
Combining this with our results for CGL
gives the following sufficient conditions.

\begin{restatable}{proposition}{pUnique}\label{prop:p-unique}
    Let \( \lambda > 0 \)
    and suppose that the convex ReLU problem
    has a unique solution.
    Then the ReLU model solution is p-unique.
    In particular,
    if \( \cbr{D_i Z \vi}_\equi \) are linearly independent,
    then the non-convex solution is p-unique.
\end{restatable}

For gated ReLU networks, it is also sufficient to check the blocks
\( \sbr{D_i Z}_{D_i \in \tilde \calD} \) to see if they satisfy GGP.
By looking at the structure of \( D_i X \),
we give simple sufficient conditions for sub-sampled
convex ReLU programs to be p-unique.

\begin{restatable}{proposition}{reluUnique}\label{prop:relu-unique}
    Let \( \lambda > 0 \) and \( p = |\tilde \calD| \).
    Suppose \( Z \) follows a continuous probability distribution
    and \( \nnz\,(D_i) \geq p \cdot d \) for every \( D_i \in \tilde \calD \).
    If \( \equi \) does not contain two blocks with the same
    activation pattern, then
    the sub-sampled convex ReLU program has a p-unique solution
    almost surely.
\end{restatable}

\cref{prop:relu-unique} requires \( n \) to be much greater than
\( d \) to be useful due to the trivial bound \( \nnz\,(D_i) \leq n \).
In practice, the condition on activation patterns can be enforced by
constraining \( v_i = 0 \) or \( w_i = 0 \) for each
activation pattern before solving the convex reformulation.

\textbf{Pruning}:
If \( (v, u) \) is a minimal solution to the convex reformulation, then the
corresponding ReLU network is the p-unique model using only those
activation patterns (Propositions~\ref{prop:minimal-solutions} and~\ref{prop:p-unique}).
Thus, \cref{alg:pruning-solutions} can be used to prune
any solution to obtain a neuron-sparse neural network
achieving the optimal training objective.
\cref{alg:pruning-solutions-nn} specializes our pruning algorithm
to the ReLU problem and extends it to support approximate pruning.
Note the resulting procedure is completely independent of the convex
reformulation.
The complexity of this method is as follows.

\begin{restatable}{proposition}{minimalComplexity}\label{prop:minimal-complexity}
    Suppose \( r = \text{rank}(X) \).
    Then an optimal and minimal ReLU network with at most
    \( m^* \leq n \) non-zero neurons can be computed in
    \( O\rbr{d^3 r^3 (n/r)^{3r}} \)
    time.
\end{restatable}

As a consequence, the complexity of computing
an optimal and minimal ReLU network is fully polynomial when \( r \)
is bounded.
We also have a more sensitive statement for the minimal width: if
\( (W^{*(1)}, W^{*(2)}) \) are optimal weights for the ReLU model, then
\( m^* \) is at most the dimension of
\( \Span \cbr{(X W^{*(1)}_{\cdot, i})_+}_i \)
and exactly this dimension under group dependence
(\cref{prop:minimal-sol-characterization}).
We experiment with pruning ReLU networks using this approach in
\cref{sec:sol-fn-experiments} and that show it
is more effective than naive pruning strategies.

\textbf{Continuity}:
First we give a negative result for singular networks, that is, models
where \( m < m^* \) and no convex reformulation exists.
In this setting, the solution map can be made to behave
arbitrarily poorly.
\begin{restatable}{proposition}{oneNeuron}\label{prop:one-neuron}
    There exists \( (Z, y) \) for which \( \calO^* \)
    is not open nor is the model fit \( f_{W^{(1)}, W^{(2)}}(Z) \) continuous in
    \( \lambda \).
\end{restatable}
%The proof for \cref{prop:one-neuron} uses an alternative reformulation of
%two-layer ReLU network which holds regardless of network width.
Combined with the next result, \cref{prop:one-neuron} indicates
that the threshold \( m^* \) may be crucial for continuity to extend
to the non-convex parameterization.
\begin{corollary}\label{cor:relu-continuity}
    Suppose \( m \geq m^* \).
    Then the optimal model fit for two-layer gated ReLU networks
    is continuous at all \( \lambda > 0 \).
    Similarly, if the (gated) ReLU solution is p-unique on an open
    interval \( \Lambda \), then the regularization path is also continuous
    on \( \Lambda \) up to permutations of the weights.
\end{corollary}
Together, \cref{cor:relu-continuity} and \cref{prop:relu-unique}
are concrete conditions for the model fit and regularization path
of a sub-sampled problem to be continuous.

\textbf{Min-Norm Solutions}:
In \cref{sec:min-norm}, we examined minimum \( \ell_2 \)-norm solutions to CGL.
However, all optimal ReLU models have the same
\( \ell_2 \)-norm when \( \lambda > 0 \).
%This follows because (a) \( p^*(\lambda) \) is
%constant over the optimal set and (b) the model fit is also constant, so that
%re-arranging the objective implies the weight-decay penalty
%is constant as well.
Minimizing the Euclidean norm of solutions to the convex reformulation instead
minimizes the sum of norms to the fourth power.
\begin{restatable}{lemma}{minNormNN}\label{lemma:min-norm-nn}
    The minimum \( \ell_2 \)-norm solution to the convex reformulation
    of a (gated) ReLU model corresponds to the
    optimal neural network which minimizes,
    \[
        r(W^{(1)}, W^{(2)}) = \sum_{i = 1}^m \norm{W^{(1)}_{\cdot, i}}_2^4 + \norm{W^{(2)}_{i}}_2^4,
    \]
    and admits no neuron-merge symmetries.
\end{restatable}

Here, neuron-merge symmetry refers to the fact that neurons with collinear
weight vectors can be merged without changing the predictions or norm of
the network (see \cref{lemma:mapping-equality} for more details).
As a result, we can compute the \( r \)-minimal optimal ReLU network
by solving Problem~\eqref{eq:min-norm-program}.
If \( \calS_\lambda \) is unknown, then using \( \act(w) \) for some solution
\( w \) as an approximation gives the \( r \)-minimal network
using a subset of those activations.

\textbf{Sensitivity}:
\cref{prop:local-sol-fn} extends similar results for the group lasso by
\citet{vaiter2012dof} to CGL using standard
CQs.
Since \( K \) is block-diagonal,
LICQ will be satisfied whenever the rows of \( Z \) are linearly independent.
SCS is more challenging; while the classical theorem of
\citet{goldman20164} establishes that SCS is satisfied for linear programs,
it is known that SCS can fail for general cone programs \citep{tunccel2012strong}.
As such, SCS must be checked on a per-problem basis in general.

In the context of gated ReLU problems, \( K = 0 \) and there is no requirement
for CSC/LICQ.
Minimal models \( w(\lambda, y) \) are weakly differentiable, which
\citet{vaiter2012dof} uses to compute the degrees of freedom of
\( w \) via Stein's Lemma \citep{stein1981estimation}.
It is straightforward to extend this calculation to the gated ReLU weights
using the chain rule, which can then be used to calculate
Stein's unbiased risk estimator.



